\section{Explainability Methods and SHAP Grouping}

The explainability layer in this fork has two jobs. First, it should make agent
behaviour inspectable for debugging and comparison. Second, it should turn
policy decisions into research artefacts that can be compared across red
profiles, reward modes, and agent families. This section therefore separates
three levels of explanation: environment-level traces, reward decomposition, and
policy-level feature attribution.

\subsection{Method Overview}

\paragraph{Environment and action traces.}
The most direct explanation method is the step log produced by
\texttt{explain.py}. It records the selected blue action, red profile, reward
components, and optional agent-provided metadata. These traces are not a
post-hoc model explanation; they are the factual execution record. They are the
right source for questions such as ``which action was taken?'', ``which red
profile was active?'', and ``when did reward loss occur?''.

\paragraph{Reward decomposition.}
Reward decomposition attributes scalar reward changes to the environment's
reward components. This is useful because the CC4 score mixes availability,
confidentiality, integrity, green activity, and phase-dependent effects. A
policy can have a good mean score while still losing reward through a specific
failure mode, for example late Restore timing or contractor-network damage.
Reward decomposition explains the consequence of decisions, whereas SHAP
explains the inputs that made a decision predictable.

\paragraph{Policy-level SHAP.}
SHAP values approximate each feature's contribution to a model output relative
to a background distribution \cite{lundberg2017unified}. In this repository SHAP
is applied to a surrogate classifier trained on observed policy decisions. The
surrogate maps logged \texttt{shap\_features} to the chosen action class. This is
model-agnostic and works for both rule-based and learned agents, but the
interpretation is intentionally narrow: the SHAP plots explain the surrogate's
approximation of the observed policy, not the true simulator dynamics and not a
causal intervention on the environment.

\subsection{Current SHAP Data Flow}

Both the v11b heuristic adapter and the MARL graph wrapper now export comparable
\texttt{shap\_features}. The current feature set contains:

\begin{itemize}
  \item previous action metadata: \texttt{prev\_action\_success}, \texttt{prev\_action};
  \item local observation flags extracted from the raw observation dictionary:
        \texttt{has\_cmd\_sh}, \texttt{has\_escalate}, \texttt{has\_decoy\_exploit};
  \item inter-agent message features:
        \texttt{msg\_any\_received}, \texttt{msg\_n\_received},
        \texttt{msg\_any\_scanned}, \texttt{msg\_any\_compromised},
        \texttt{msg\_n\_scanned\_flags}, and
        \texttt{msg\_n\_compromised\_flags}.
\end{itemize}

The current implementation is useful as a first comparable pipeline, but it is
coarse. Keyword-derived features such as \texttt{has\_cmd\_sh} or
\texttt{has\_escalate} compress many hosts and subnets into one bit. This makes
the SHAP plots stable and cheap, but it can hide where the signal came from.

\subsection{Recommended Semantic Feature Groups}

If the SHAP pipeline were restarted from the current observation interface, the
feature schema should be grouped around operational cyber-defence concepts
rather than around parser details. The recommended groups are:

\begin{longtable}{p{0.24\textwidth}p{0.30\textwidth}p{0.36\textwidth}}
\toprule
\textbf{Group} & \textbf{Example features} & \textbf{Interpretation} \\
\midrule
\texttt{action\_state} &
\texttt{prev\_action\_success}, \texttt{prev\_action\_type},
\texttt{in\_progress\_flag} &
Whether the agent is reacting to the outcome or duration of its previous action. \\
\texttt{red\_presence} &
\texttt{n\_cmd\_sh}, \texttt{n\_privileged\_sessions},
\texttt{n\_malware\_files}, \texttt{red\_seen\_any} &
Evidence that red is present on observed hosts. \\
\texttt{red\_intent} &
\texttt{n\_escalation\_signals}, \texttt{n\_impact\_adjacent\_signals},
\texttt{n\_degrade\_precursors} &
Evidence about what red is likely trying to do next. \\
\texttt{asset\_risk} &
\texttt{host\_is\_server}, \texttt{subnet\_criticality},
\texttt{contractor\_subnet\_flag}, \texttt{phase} &
How expensive it would be to leave the observed state untreated. \\
\texttt{communication} &
\texttt{msg\_n\_received}, \texttt{msg\_n\_scanned\_flags},
\texttt{msg\_n\_compromised\_flags}, \texttt{msg\_sender\_count} &
How much of the decision is driven by shared information from other blue agents. \\
\texttt{decoy\_surface} &
\texttt{n\_active\_decoys}, \texttt{n\_decoy\_hits},
\texttt{decoy\_depleted\_flag}, \texttt{red\_discover\_deception\_seen} &
How deception-related mechanics shape blue's decision. This name is preferred
over the broader label \texttt{deception}. \\
\texttt{network\_control} &
\texttt{blocked\_subnet\_count}, \texttt{comms\_blocked\_flag},
\texttt{known\_reachable\_subnets} &
Whether the agent is acting because of connectivity or containment constraints. \\
\texttt{uncertainty} &
\texttt{unknown\_host\_count}, \texttt{stale\_observation\_age},
\texttt{conflicting\_signals} &
How much the action is driven by partial observability rather than direct evidence. \\
\bottomrule
\end{longtable}

The important change is that \texttt{decoy\_surface} should not be treated as a
generic explanation bucket. In CC4, deception has a concrete simulator meaning:
blue can deploy deceptive services through \texttt{Misinform}, and red can spend
time on \texttt{DiscoverDeception} to detect decoys. A SHAP group with that name
should therefore contain only decoy deployment, decoy hit/depletion, and
red-deception-check signals. If those signals are not available in the blue
observation, the group should be omitted or marked as environment-side metadata,
not inferred through unrelated strings.

\subsection{How to Improve SHAP Quality}

The next iteration should improve the SHAP values in four ways.

\paragraph{Use structured feature extraction.}
The current keyword scan is robust to changing dictionary shapes, but it is not
semantically precise. The extractor should parse the observation dictionary into
stable typed counters: per-host compromise indicators, process flags, malware
file flags, service observations, message flags, subnet identifiers, and mission
phase. This reduces accidental feature matches and makes the feature names
defensible in a paper.

\paragraph{Keep host-local and aggregate views.}
A single global flag such as \texttt{has\_cmd\_sh} answers whether a signal
exists, not where it exists or how many hosts are affected. The better schema is
to log both aggregate counters and local action-target features. For example,
\texttt{n\_cmd\_sh\_observed} describes the agent's whole observation, while
\texttt{target\_has\_cmd\_sh} describes the host selected by the action. This is
especially important for distinguishing \texttt{Analyse}, \texttt{Remove},
\texttt{Restore}, and \texttt{Misinform}.

\paragraph{Explain action families separately.}
The current multiclass surrogate compresses all action classes into one
classifier. For a paper figure, it is often clearer to train one-vs-rest
surrogates for action families such as \texttt{Monitor}, \texttt{Analyse},
\texttt{Remove}, \texttt{Restore}, \texttt{Misinform}, and \texttt{Block}. This
turns the explanation into a question readers can understand: ``what evidence
makes Restore more likely than other actions?''

\paragraph{Use matched background samples.}
SHAP values depend on the background distribution. For CC4, background rows
should be stratified by red profile, mission phase, and agent index. Otherwise
the explanation may mostly rediscover that some profiles or phases are more
common in the sample. Large runs should continue to evaluate all episodes while
using deterministic SHAP sampling through \texttt{--shap-episode-stride},
\texttt{--shap-step-stride}, and \texttt{--shap-max-rows-per-profile}.

\subsection{Recommended Near-Term Implementation}

The practical next step is not to add more plots first. It is to make the
\texttt{shap\_features} contract stronger:

\begin{enumerate}
  \item replace the keyword-only extractor with a structured observation
        extractor that emits typed counters and target-local features;
  \item rename the current broad \texttt{deception} group to
        \texttt{decoy\_surface};
  \item store a \texttt{feature\_group} mapping next to each SHAP CSV so grouped
        plots are reproducible;
  \item report both per-feature mean absolute SHAP and grouped mean absolute
        SHAP, normalized within each action family;
  \item keep one compact global summary table, but use per-action-family plots
        for interpretation in the paper.
\end{enumerate}

This produces less flashy but more defensible explanations: the plots would show
whether the agent's \texttt{Restore} decisions are driven by red-presence
signals, asset risk, communication, or uncertainty, and whether this changes
under different red profiles.
